% !TeX root = ../教材题库_全部条目_转写工作稿.tex
% Source: XJ-B1-C05-TH-P173-Y0178
% 第5章 / 5.2 任意角的三角函数 / 教材 p.173、174 / PDF 第180、181页
% OCR未校对：公式与文字均可能有误；未完成数学TeX转写；禁止直接发给学生
% 原题及相邻原解截图（核对用）：../题目原图/XJ-B1-C05-TH-P173-Y0178-01.webp, ../题目原图/XJ-B1-C05-TH-P173-Y0178-02.webp
% 本文件只在转写完成后作为题干使用；OCR 条目尚不是可靠题干。
\begin{bankitem}{p.173 思考（含紧邻说明）}{第5章 / 5.2 任意角的三角函数 / 教材 p.173、174 / PDF 第180、181页}{XJ-B1-C05-TH-P173-Y0178}{OCR未校对：公式与文字均可能有误；未完成数学TeX转写；禁止直接发给学生}
\begin{ocrdraft}
[教材p.173 / PDF第180页]
\par
？思考在初中学习锐角三角函数时，根据直角三角形两锐角互余的关系得出
\par
了锐角仅与它的余角丌
\par
的三角函数之间的关系：
\par
2
\par
Sln
\par
．一COS，COS
\par
Sln\ensuremath{\cdot}
\par
2
\par
上述关系式是否对任意角成立呢？
\par
如图5．2一10，设角仅的终边OM交单位圆于点P(x，回，
\par
角《一。的终边OMI交单位圆于点PI.
\par
由于角。与角《一“的弧度数的平均值为砉[。十《一。
\par
图5．2一10
\par
，于是角的终边OM与角丌
\par
仅的终边OMI关于Z主ON一
\par
4
\par
匹的终边ON所在直线3-对称，因此点P(x，y)与PI关于直线3，一对称，并且
\par
4
\par
点PI的坐标为巧，，x).
\par
由三角函数的定义可知：
\par
S1n仅．一'
\par
COS
\par
2
\par
-一30=S1n仅，
\par
又利用公式二可得：
\par
S1n一一\ocrglyph{00AF}\ocrglyph{00E5}\ocrglyph{00AF}a
\par
COS仅，
\par
\ocrglyph{00AF}sin(—a)
\par
cos
\par
一cos
\par
S1na\ensuremath{\cdot}
\par
由此得出关于角\ensuremath{\pm}与的正、余弦关系式：
\par
公式五
\par
COS，
\par
COS
\par
仅
\par
Slna，
\par
你能用单位圆中的三
\par
2
\par
角函数线推导公式五吗？
\par
COS，
\par
Sln\ensuremath{\cdot}
\par
公式五可概括为如下法则：
\par
匹\ensuremath{\pm}a的正弦（余弦）函数值，等于角a的余弦（正弦）函数值，前面添上一个把角
\par
2
\par
仅看成锐角时原函数值的符号．
\par
利用公式五，可以实现正弦函数与余弦函数的相互转化．
\par
[教材p.174 / PDF第181页]
\par
当角的终边不在坐标轴上时，还可以得出以下公式：
\par
公式六
\par
Sln
\par
COSa
\par
tan
\par
Sln仅
\par
tan
\par
COS
\par
Sln
\par
玉十a
\par
COS仅
\par
tan
\par
tana
\par
cos一十仅
\par
Sln
\par
以上关于角a与2差丌十a（差ez），
\par
匹\ensuremath{\pm}仅的三角
\par
丌
\par
仅，
\par
函数的关系式，统称诱导公式．
\par
诱导公式的一个重要
\par
作用是可以将不是锐角的
\end{ocrdraft}
\end{bankitem}
